\subsubsection{Receipts and write-offs between annual dates}
\label{sec:dated-credit-schedules}

The terminal claim isolates how an initial credit loss differs from a later
change in that loss. A bank also needs to know what happens when cash arrives
in instalments and a portion must be written off before the next year end.
Consider a separate purchased credit-impaired asset acquired on 1 January 2025.
The price is 98 and an integral transaction cost is 2, giving initial carrying
amount 100. Contractual receipts are 66 on 1 January 2026 and 72.6 on
1 January 2027. Initial expected receipts are 55 and 60.5. Assume the same
separately established POCI and amortised-cost classifications as before, and
no modification, foreign exchange or tax. Opening cash and capital are 300.

The credit-adjusted effective rate discounts the expected receipts to the
amount invested, including the integral cost. Here it is 10\% because
\[
 \frac{55}{1.1}+\frac{60.5}{1.1^2}=50+50=100.
\]
The corresponding contractual present value is 120. Its excess of 20 over
the expected value is embedded in the acquisition measurement; the subsequent-
change allowance again starts at zero. The Appendix A definition and
B5.4.1--4 require the fee and cost classification to precede rate calculation
\citep[IFRS~9, Appendix~A and B5.4.1--4]{eu20231803}. Calling the 2 a fee
does not, by itself, establish that it is integral.

For intermediate dates, this example supplies ACT/365F with annual compound
interest: count actual elapsed days and divide by 365. This is a declared
calculation convention, rather than a day-count rule prescribed by IFRS~9.
Let \(F(s,t)\) be the accumulation factor from date \(s\) to date \(t\),
and retain \(G\), \(L\) and \(N=G-L\) for gross carrying, signed
allowance and net carrying. During an interval without cash or new credit
information, both balances accrete at the original rate:
\[
 F(s,t)=(1+r)^{\operatorname{days}(s,t)/365},\qquad
 G_t=G_sF(s,t),\quad L_t=L_sF(s,t).
\]
Their difference therefore earns interest
\(N_s\{F(s,t)-1\}\). At a credit reassessment, the bank discounts the
complete remaining expected receipts \(e_d\), each at its own date \(d\):
\[
 N_t=\sum_{d\geq t}\frac{e_d}{F(t,d)},\qquad L_t=G_t-N_t.
\]
The change in \(N_t\) due to revised expectations is impairment income or
expense; the preceding passage-of-time change is interest. This distinction
implements the original-rate requirement and subsequent-loss treatment in
5.4.1(a), 5.5.13--14 and B5.5.45 \citep{eu20231803}.

There are 181 days from acquisition to 1 July 2025. Write
\(f=(11/10)^{181/365}\). Immediately before new information, gross and
net carrying are both \(100f\). Suppose expected receipts fall to 44 and
48.4. Each is four-fifths of its initial expectation, so their combined
present value is \(80f\). The bank recognizes an impairment loss of
\(20f\), leaving \(G=100f\), \(L=20f\) and \(N=80f\).

Now add a separate finding: an identified, proportional 20\% portion of the
whole claim has no reasonable prospect of recovery. Its expected cash is zero;
the stated allocation assigns it 20\% of gross carrying. The financial
write-off is consequently \(W=20f\). Reducing both gross and allowance by
\(W\) leaves gross \(80f\), allowance zero and net \(80f\). There is
no second impairment charge. Paragraphs 5.4.4 and B5.4.9 require this earlier
write-off when that finding applies, even though the surviving claim still
has positive expected receipts \citep{eu20231803}. The proportional allocation
is an explicit assumption of this example; another identified portion needs
its own supported allocation.

The removed contractual receipts are 13.2 and 14.52, leaving 52.8 and 58.08
in the financial claim. The bank retains the removed amounts and leaves
legal forgiveness unassessed. This distinction matters when a later recovery
or legal-enforcement question arises. A lower expected total alone would not
have supplied the no-recovery finding or identified the portion to remove.

At 1 January 2026, the remaining gross balance has grown to 88. Collection
of 44 leaves gross and net carrying of 44. Immediately afterward, a supported
improvement raises the final expected receipt to 55, within the surviving
contractual 58.08. Its present value is \(55/1.1=50\), so the bank records
an impairment gain of 6 and a signed allowance of \(44-50=-6\).
The following year's net interest is 5. Just before the final collection,
gross is 48.4, allowance is \(-6.6\), and net carrying is 55. Receiving
55 and removing gross and allowance together closes the asset without another
gain. Final cash is
\[
 300-98-2+44+55=299.
\]
The lifetime loss is therefore 1. The period origins give the same answer:
cumulative interest is \(20f-7\) and net impairment is \(-20f+6\),
whose sum is \(-1\). The cancellation does not make either interim amount
an exact rational number.

A performing asset provides a useful comparison. Suppose the bank instead
advances 102 and receives an integral origination fee of 2, with an independent
zero-credit-loss assessment. The same expected contractual receipts 55 and
60.5 imply the same initial carrying amount 100 and effective rate 10\%.
On 1 July, a revised estimate under an existing prepayment right changes the
remaining receipt to 105 on the following 1 January. Assume this is a
non-credit estimate change, not a contractual modification. Under B5.4.6,
the bank recomputes gross carrying at the original rate:
\[
 G_{\rm new}=\frac{105}{1.1/f}=\frac{1050}{11}f,
 \qquad
 G_{\rm new}-100f=-\frac{50}{11}f.
\]
The negative catch-up enters profit separately from credit impairment.
Subsequent interest brings the asset to 105 for collection. Total profit is
\(2-102+105=5\). Omitting the integral fee would fit a different rate to
102 and would misstate this comparison.

Dates also change the numerical problem. Fractional powers such as \(f\)
are generally irrational, so a rational calculator must preserve a bound
rather than label a rounded decimal exact. Reduce the day fraction to
\(p/q\) and choose rational endpoints \(a\) and \(b\) on a declared
decimal grid. The following integer-power comparisons establish containment:
\[
 a^q\leq(1+r)^p\leq b^q
 \quad\Longrightarrow\quad
 a\leq F(s,t)\leq b.
\]
All quantities are positive, which justifies taking the monotone \(q\)th
root. Interval arithmetic then carries these bounds through discounting and
each balance; the requested scale and accumulated widths remain part of the
result. A leap-day interval uses its actual number of days, rather than an
implicit annual-anniversary convention. Zero and negative rates above
\(-1\) are admitted under the same convention.

For positive future receipts, present value decreases continuously as the rate
increases. Certified endpoint signs can therefore enclose a unique effective
rate. Such a bracket is useful numerical evidence, but its midpoint is not
an exact rate. The two financial examples require an exactly established
initial PV identity and use bounds for the intervening dates. Propagating an
uncertain rate through a general financial schedule remains a separate method.

The dated records retain cash, write-off and catch-up origins so gross carrying
can be reconstructed from their independently compounded contributions.
Terminal allowance clearing is also an origin; omitting it would leave the
gross reconciliation wrong when the allowance is negative. Every event needs
its information date, unchanged-terms assessment and any required write-off
finding. A missing receipt or unsupported finding returns no partial report.
These finite methods supply dated accounting consequences; credit forecasting,
arbitrary write-off allocation and the wider bank-company joins remain open.
